\IfFileExists{stacks-project.cls}{%
\documentclass{stacks-project}
}{%
\documentclass{amsart}
}

% For dealing with references we use the comment environment
\usepackage{verbatim}
\newenvironment{reference}{\comment}{\endcomment}
%\newenvironment{reference}{}{}
\newenvironment{slogan}{\comment}{\endcomment}
\newenvironment{history}{\comment}{\endcomment}

% For commutative diagrams we use Xy-pic
\usepackage[all]{xy}

% We use 2cell for 2-commutative diagrams.
\xyoption{2cell}
\UseAllTwocells

% We use multicol for the list of chapters between chapters
\usepackage{multicol}

% This is generally recommended for better output
\usepackage{lmodern}
\usepackage[T1]{fontenc}

% For cross-file-references
\usepackage{xr-hyper}

% Package for hypertext links:
\usepackage{hyperref}

% For any local file, say "hello.tex" you want to link to please
% use \externaldocument[hello-]{hello}
\externaldocument[introduction-]{introduction}
\externaldocument[conventions-]{conventions}
\externaldocument[sets-]{sets}
\externaldocument[categories-]{categories}
\externaldocument[topology-]{topology}
\externaldocument[sheaves-]{sheaves}
\externaldocument[sites-]{sites}
\externaldocument[stacks-]{stacks}
\externaldocument[fields-]{fields}
\externaldocument[algebra-]{algebra}
\externaldocument[brauer-]{brauer}
\externaldocument[homology-]{homology}
\externaldocument[derived-]{derived}
\externaldocument[simplicial-]{simplicial}
\externaldocument[more-algebra-]{more-algebra}
\externaldocument[smoothing-]{smoothing}
\externaldocument[modules-]{modules}
\externaldocument[sites-modules-]{sites-modules}
\externaldocument[injectives-]{injectives}
\externaldocument[cohomology-]{cohomology}
\externaldocument[sites-cohomology-]{sites-cohomology}
\externaldocument[dga-]{dga}
\externaldocument[dpa-]{dpa}
\externaldocument[sdga-]{sdga}
\externaldocument[hypercovering-]{hypercovering}
\externaldocument[schemes-]{schemes}
\externaldocument[constructions-]{constructions}
\externaldocument[properties-]{properties}
\externaldocument[morphisms-]{morphisms}
\externaldocument[coherent-]{coherent}
\externaldocument[divisors-]{divisors}
\externaldocument[limits-]{limits}
\externaldocument[varieties-]{varieties}
\externaldocument[topologies-]{topologies}
\externaldocument[descent-]{descent}
\externaldocument[perfect-]{perfect}
\externaldocument[more-morphisms-]{more-morphisms}
\externaldocument[flat-]{flat}
\externaldocument[groupoids-]{groupoids}
\externaldocument[more-groupoids-]{more-groupoids}
\externaldocument[etale-]{etale}
\externaldocument[chow-]{chow}
\externaldocument[intersection-]{intersection}
\externaldocument[pic-]{pic}
\externaldocument[weil-]{weil}
\externaldocument[adequate-]{adequate}
\externaldocument[dualizing-]{dualizing}
\externaldocument[duality-]{duality}
\externaldocument[discriminant-]{discriminant}
\externaldocument[derham-]{derham}
\externaldocument[local-cohomology-]{local-cohomology}
\externaldocument[algebraization-]{algebraization}
\externaldocument[curves-]{curves}
\externaldocument[resolve-]{resolve}
\externaldocument[models-]{models}
\externaldocument[functors-]{functors}
\externaldocument[equiv-]{equiv}
\externaldocument[pione-]{pione}
\externaldocument[etale-cohomology-]{etale-cohomology}
\externaldocument[proetale-]{proetale}
\externaldocument[relative-cycles-]{relative-cycles}
\externaldocument[more-etale-]{more-etale}
\externaldocument[trace-]{trace}
\externaldocument[crystalline-]{crystalline}
\externaldocument[spaces-]{spaces}
\externaldocument[spaces-properties-]{spaces-properties}
\externaldocument[spaces-morphisms-]{spaces-morphisms}
\externaldocument[decent-spaces-]{decent-spaces}
\externaldocument[spaces-cohomology-]{spaces-cohomology}
\externaldocument[spaces-limits-]{spaces-limits}
\externaldocument[spaces-divisors-]{spaces-divisors}
\externaldocument[spaces-over-fields-]{spaces-over-fields}
\externaldocument[spaces-topologies-]{spaces-topologies}
\externaldocument[spaces-descent-]{spaces-descent}
\externaldocument[spaces-perfect-]{spaces-perfect}
\externaldocument[spaces-more-morphisms-]{spaces-more-morphisms}
\externaldocument[spaces-flat-]{spaces-flat}
\externaldocument[spaces-groupoids-]{spaces-groupoids}
\externaldocument[spaces-more-groupoids-]{spaces-more-groupoids}
\externaldocument[bootstrap-]{bootstrap}
\externaldocument[spaces-pushouts-]{spaces-pushouts}
\externaldocument[spaces-chow-]{spaces-chow}
\externaldocument[groupoids-quotients-]{groupoids-quotients}
\externaldocument[spaces-more-cohomology-]{spaces-more-cohomology}
\externaldocument[spaces-simplicial-]{spaces-simplicial}
\externaldocument[spaces-duality-]{spaces-duality}
\externaldocument[formal-spaces-]{formal-spaces}
\externaldocument[restricted-]{restricted}
\externaldocument[spaces-resolve-]{spaces-resolve}
\externaldocument[formal-defos-]{formal-defos}
\externaldocument[defos-]{defos}
\externaldocument[cotangent-]{cotangent}
\externaldocument[examples-defos-]{examples-defos}
\externaldocument[algebraic-]{algebraic}
\externaldocument[examples-stacks-]{examples-stacks}
\externaldocument[stacks-sheaves-]{stacks-sheaves}
\externaldocument[criteria-]{criteria}
\externaldocument[artin-]{artin}
\externaldocument[quot-]{quot}
\externaldocument[stacks-properties-]{stacks-properties}
\externaldocument[stacks-morphisms-]{stacks-morphisms}
\externaldocument[stacks-limits-]{stacks-limits}
\externaldocument[stacks-cohomology-]{stacks-cohomology}
\externaldocument[stacks-perfect-]{stacks-perfect}
\externaldocument[stacks-introduction-]{stacks-introduction}
\externaldocument[stacks-more-morphisms-]{stacks-more-morphisms}
\externaldocument[stacks-geometry-]{stacks-geometry}
\externaldocument[moduli-]{moduli}
\externaldocument[moduli-curves-]{moduli-curves}
\externaldocument[examples-]{examples}
\externaldocument[exercises-]{exercises}
\externaldocument[guide-]{guide}
\externaldocument[desirables-]{desirables}
\externaldocument[coding-]{coding}
\externaldocument[obsolete-]{obsolete}
\externaldocument[fdl-]{fdl}
\externaldocument[index-]{index}

% Theorem environments.
%
\theoremstyle{plain}
\newtheorem{theorem}[subsection]{Theorem}
\newtheorem{proposition}[subsection]{Proposition}
\newtheorem{lemma}[subsection]{Lemma}

\theoremstyle{definition}
\newtheorem{definition}[subsection]{Definition}
\newtheorem{example}[subsection]{Example}
\newtheorem{exercise}[subsection]{Exercise}
\newtheorem{situation}[subsection]{Situation}

\theoremstyle{remark}
\newtheorem{remark}[subsection]{Remark}
\newtheorem{remarks}[subsection]{Remarks}

\numberwithin{equation}{subsection}

% Macros
%
\def\lim{\mathop{\mathrm{lim}}\nolimits}
\def\colim{\mathop{\mathrm{colim}}\nolimits}
\def\Spec{\mathop{\mathrm{Spec}}}
\def\Hom{\mathop{\mathrm{Hom}}\nolimits}
\def\Ext{\mathop{\mathrm{Ext}}\nolimits}
\def\SheafHom{\mathop{\mathcal{H}\!\mathit{om}}\nolimits}
\def\SheafExt{\mathop{\mathcal{E}\!\mathit{xt}}\nolimits}
\def\Sch{\mathit{Sch}}
\def\Mor{\mathop{\mathrm{Mor}}\nolimits}
\def\Ob{\mathop{\mathrm{Ob}}\nolimits}
\def\Sh{\mathop{\mathit{Sh}}\nolimits}
\def\NL{\mathop{N\!L}\nolimits}
\def\CH{\mathop{\mathrm{CH}}\nolimits}
\def\proetale{{pro\text{-}\acute{e}tale}}
\def\etale{{\acute{e}tale}}
\def\QCoh{\mathit{QCoh}}
\def\Ker{\mathop{\mathrm{Ker}}}
\def\Im{\mathop{\mathrm{Im}}}
\def\Coker{\mathop{\mathrm{Coker}}}
\def\Coim{\mathop{\mathrm{Coim}}}

% Boxtimes
%
\DeclareMathSymbol{\boxtimes}{\mathbin}{AMSa}{"02}

%
% Macros for moduli stacks/spaces
%
\def\QCohstack{\mathcal{QC}\!\mathit{oh}}
\def\Cohstack{\mathcal{C}\!\mathit{oh}}
\def\Spacesstack{\mathcal{S}\!\mathit{paces}}
\def\Quotfunctor{\mathrm{Quot}}
\def\Hilbfunctor{\mathrm{Hilb}}
\def\Curvesstack{\mathcal{C}\!\mathit{urves}}
\def\Polarizedstack{\mathcal{P}\!\mathit{olarized}}
\def\Complexesstack{\mathcal{C}\!\mathit{omplexes}}
% \Pic is the operator that assigns to X its picard group, usage \Pic(X)
% \Picardstack_{X/B} denotes the Picard stack of X over B
% \Picardfunctor_{X/B} denotes the Picard functor of X over B
\def\Pic{\mathop{\mathrm{Pic}}\nolimits}
\def\Picardstack{\mathcal{P}\!\mathit{ic}}
\def\Picardfunctor{\mathrm{Pic}}
\def\Deformationcategory{\mathcal{D}\!\mathit{ef}}
\title{The normalised complex functor is faithful and reflects isomorphisms, injections and surjections}
\author{The Stacks Project Authors}
\date{stacks-project:simplicial.tex:local-snapshot-20261012; GNU FDL 1.2 or later}
\begin{document}
\maketitle
\noindent\textit{Editorial note.} The counted statement is the author's own lemma environment labelled lemma-N-faithful (source lines 4216-4221); the proof environment begins on the second line after it and is its own proof, closing at source line 4229.\par\medskip
\noindent\textit{Editorial note.} The proof invokes two author results of the same chapter by label; both are retained as context spans and the generator expands the context transitively (budget 800 lines for this card), while reference targets in other chapters that carry canonical project tags are auto-declared as non-load-bearing source pointers.\par\medskip
\noindent\textit{Editorial note.} The retained preamble is the chapter's own minimal preamble (sources/stacks-simplicial/preamble.tex, taken with its real line count), a verbatim excerpt of the native Xy-pic setup, the native theorem declarations with the mathematical macros, and the project's own no-op reference environments; no macro was invented or redefined.\par\medskip
\noindent\textit{Editorial scope.} One original result of the Stacks Project's Simplicial Methods chapter, counted once: the normalised complex functor is faithful and reflects isomorphisms, injections and surjections.\par\medskip
\noindent\textit{Allowed standard prerequisites (editorial).}\begin{itemize}
\item The normalised complex functor N. The functor N from simplicial objects of an abelian category to chain complexes, with the canonical comparison map from N to s.. It is the functor of the statement (source lines 4218-4220).
\item The canonical splitting of the comparison map. The comparison map N(U) to s(U) is split, so that a morphism is determined by its effect on the normalised subobjects (the author's lemma-splitting-abelian-category, retained as context).. It gives faithfulness (source lines 4224-4225).
\item The injective-map lemma. The criterion identifying the maps of simplicial abelian objects whose normalised complex map is injective (the author's lemma-injective-map-simplicial-abelian, retained as context).. It gives the reflection of injections, surjections and isomorphisms (source lines 4226-4228).
\end{itemize}\medskip
\begin{lemma}
\label{lemma-splitting-abelian-category}
Let $\mathcal{A}$ be an abelian category.
Let $U$ be a simplicial object in $\mathcal{A}$.
Then $U$ has a splitting obtained by taking $N(U_0) = U_0$ and
for $m \geq 1$ taking
$$
N(U_m) = \bigcap\nolimits_{i = 0}^{m - 1} \Ker(d^m_i).
$$
Moreover, this splitting is functorial on the category of
simplicial objects of $\mathcal{A}$.
\end{lemma}

\begin{proof}
For any object $A$ of $\mathcal{A}$ we obtain
a simplicial abelian group $\Mor_\mathcal{A}(A, U)$.
Each of these are canonically split by Lemma
\ref{lemma-splitting-simplicial-groups}. Moreover,
$$
N(\Mor_\mathcal{A}(A, U_m)) =
\bigcap\nolimits_{i = 0}^{m - 1} \Ker(d^m_i) =
\Mor_\mathcal{A}(A, N(U_m)).
$$
Hence we see that the morphism (\ref{equation-splitting})
becomes an isomorphism after applying the functor
$\Mor_\mathcal{A}(A, -)$ for any object of $\mathcal{A}$.
Hence it is an isomorphism by the Yoneda lemma.
\end{proof}

\begin{lemma}
\label{lemma-injective-map-simplicial-abelian}
\begin{slogan}
The Dold-Kan normalization functor reflects
injectivity, surjectivity, and isomorphy.
\end{slogan}
Let $\mathcal{A}$ be an abelian category.
Let $f : U \to V$ be a morphism of
simplicial objects of $\mathcal{A}$.
If the induced morphisms $N(f)_i : N(U)_i \to N(V)_i$
are injective for all $i$, then $f_i$ is
injective for all $i$. Same holds with ``injective'' replaced
with ``surjective'', or ``isomorphism''.
\end{lemma}

\begin{proof}
This is clear from Lemma \ref{lemma-splitting-abelian-category}
and the definition of a splitting.
\end{proof}

\begin{equation}
\label{equation-splitting}
\coprod\nolimits_{\varphi : [n] \to [m]\text{ surjective}}
N(U_m)
\longrightarrow
U_n
\end{equation}

\begin{lemma}
\label{lemma-splitting-simplicial-groups}
Let $U$ be a simplicial abelian group.
Then $U$ has a splitting obtained by taking $N(U_0) = U_0$ and
for $m \geq 1$ taking
$$
N(U_m) = \bigcap\nolimits_{i = 0}^{m - 1} \Ker(d^m_i).
$$
Moreover, this splitting is functorial on the category
of simplicial abelian groups.
\end{lemma}

\begin{proof}
By induction on $n$ we will show that the choice of $N(U_m)$
in the lemma guarantees that (\ref{equation-splitting}) is
an isomorphism for $m \leq n$. This is clear for $n = 0$.
In the rest of this proof we are going to
drop the superscripts from the maps $d_i$ and $s_i$ in order
to improve readability. We will also repeatedly use the relations
from Remark \ref{remark-relations}.

\medskip\noindent
First we make a general remark.
For $0 \leq i \leq m$ and $z \in U_m$ we have
$d_i(s_i(z)) = z$. Hence we can write
any $x \in U_{m + 1}$ uniquely as
$x = x' + x''$ with $d_i(x') = 0$
and $x'' \in \Im(s_i)$
by taking $x' = (x - s_i(d_i(x)))$ and
$x'' = s_i(d_i(x))$. Moreover, the element
$z \in U_m$ such that $x'' = s_i(z)$
is unique because $s_i$ is injective.

\medskip\noindent
Here is a procedure for decomposing any $x \in U_{n + 1}$.
First, write $x = x_0 + s_0(z_0)$ with $d_0(x_0) = 0$.
Next, write $x_0 = x_1 + s_1(z_1)$ with
$d_1(x_1) = 0$. Continue like this to get
\begin{eqnarray*}
x & = & x_0 + s_0(z_0), \\
x_0 & = & x_1 + s_1(z_1), \\
x_1 & = & x_2 + s_2(z_2), \\
\ldots & \ldots & \ldots \\
x_{n - 1} & = & x_n + s_n(z_n)
\end{eqnarray*}
where $d_i(x_i) = 0$ for all $i = n, \ldots, 0$.
By our general remark above all of the $x_i$
and $z_i$ are determined uniquely by $x$.
We claim that
$x_i \in
\Ker(d_0) \cap
\Ker(d_1) \cap
\ldots \cap
\Ker(d_i)$
and
$z_i \in
\Ker(d_0) \cap
\ldots \cap
\Ker(d_{i - 1})$
for $i = n, \ldots, 0$.
Here and in the following
an empty intersection of kernels indicates
the whole space; i.e.,
the notation
$z_0 \in \Ker(d_0) \cap
\ldots \cap
\Ker(d_{i - 1})$
when $i = 0$ means $z_0 \in U_n$ with no restriction.

\medskip\noindent
We prove this by ascending induction on $i$.
It is clear for $i = 0$ by construction of $x_0$ and $z_0$.
Let us prove it for $0 < i \leq n$ assuming the result for $i - 1$.
First of all we have $d_i(x_i) = 0$ by construction.
So pick a $j$ with $0 \leq j < i$. We have
$d_j(x_{i - 1}) = 0$ by induction. Hence
$$
0 = d_j(x_{i - 1})
= d_j(x_i) + d_j(s_i(z_i))
= d_j(x_i) + s_{i - 1}(d_j(z_i)).
$$
The last equality by the relations of Remark \ref{remark-relations}.
These relations also imply that
$d_{i - 1}(d_j(x_i)) = d_j(d_i(x_i)) = 0$
because $d_i(x_i)= 0$ by construction.
Then the uniqueness in the general remark above shows the equality
$0 = x' + x'' = d_j(x_i) + s_{i - 1}(d_j(z_i))$
can only hold if both terms are zero. We conclude that
$d_j(x_i) = 0$ and by injectivity of $s_{i - 1}$ we also
conclude that $d_j(z_i) = 0$. This proves the claim.

\medskip\noindent
The claim implies we can uniquely write
$$
x = s_0(z_0) + s_1(z_1) + \ldots + s_n(z_n) + x_n
$$
with $x_n \in N(U_{n + 1})$ and
$z_i \in \Ker(d_0) \cap \ldots \cap \Ker(d_{i - 1})$.
We can reformulate this as saying that we have found a direct
sum decomposition
$$
U_{n + 1}
=
N(U_{n + 1})
\oplus
\bigoplus\nolimits_{i = 0}^{i = n}
s_i\Big(\Ker(d_0) \cap \ldots \cap \Ker(d_{i - 1})\Big)
$$
with the property that
$$
\Ker(d_0) \cap \ldots \cap \Ker(d_j)
=
N(U_{n + 1}) \oplus
\bigoplus\nolimits_{i = j + 1}^{i = n}
s_i\Big(\Ker(d_0) \cap \ldots \cap \Ker(d_{i - 1})\Big)
$$
for $j = 0, \ldots, n$.
The result follows from this statement as follows.
Each of the $z_i$ in the expression for $x$
can be written uniquely as
$$
z_i = s_i(z'_{i, i}) + \ldots + s_{n - 1}(z'_{i, n - 1}) + z_{i, 0}
$$
with $z_{i, 0} \in N(U_n)$ and
$z'_{i, j} \in \Ker(d_0) \cap \ldots \cap \Ker(d_{j - 1})$.
The first few steps in the decomposition of $z_i$ are zero because
$z_i$ already is in the kernel of $d_0, \ldots, d_i$.
This in turn uniquely gives
$$
x = x_n + s_0(z_{0, 0}) + s_1(z_{1, 0}) + \ldots + s_n(z_{n, 0}) +
\sum\nolimits_{0 \leq i \leq j \leq n - 1} s_i(s_j(z'_{i, j})).
$$
Continuing in this fashion we see that we in the end obtain
a decomposition of $x$ as a sum of terms
of the form
$$
s_{i_1} s_{i_2} \ldots s_{i_k} (z)
$$
with $0 \leq i_1 \leq i_2 \leq \ldots \leq i_k \leq n - k + 1$ and
$z \in N(U_{n + 1 - k})$. This is exactly the required
decomposition, because any surjective map $[n + 1] \to [n + 1 - k]$
can be uniquely expressed in the form
$$
\sigma^{n - k}_{i_k} \ldots \sigma^{n - 1}_{i_2} \sigma^n_{i_1}
$$
with $0 \leq i_1 \leq i_2 \leq \ldots \leq i_k \leq n - k + 1$.
\end{proof}

\begin{remark}
\label{remark-relations}
By abuse of notation we sometimes write $d_i : U_n \to U_{n - 1}$
instead of $d^n_i$, and similarly for $s_i : U_n \to U_{n + 1}$.
The relations among the morphisms $d^n_i$ and $s^n_i$
may be expressed as follows:
\begin{enumerate}
\item If $i < j$, then $d_i \circ d_j = d_{j - 1} \circ d_i$.
\item If $i < j$, then $d_i \circ s_j = s_{j - 1} \circ d_i$.
\item We have $\text{id} = d_j \circ s_j = d_{j + 1} \circ s_j$.
\item If $i > j + 1$, then $d_i \circ s_j = s_j \circ d_{i - 1}$.
\item If $i \leq j$, then $s_i \circ s_j = s_{j + 1} \circ s_i$.
\end{enumerate}
This means that whenever the compositions on both the left and the
right are defined then the corresponding equality should hold.
\end{remark}

\begin{lemma}
\label{lemma-N-faithful}
Let $\mathcal{A}$ be an abelian category.
The functor $N$ is faithful, and reflects
isomorphisms, injections and surjections.
\end{lemma}

\begin{proof}
The faithfulness is immediate from the canonical
splitting of Lemma \ref{lemma-splitting-abelian-category}.
The statement on reflecting injections, surjections, and
isomorphisms follows from
Lemma \ref{lemma-injective-map-simplicial-abelian}.
\end{proof}

\end{document}
